\documentclass[11pt]{article}
\usepackage{amsmath,amssymb,amsthm}
\newtheorem{proposition}{Proposition}
\title{A Product Comparison Need Not Have a Common Tangent Decoder}
\author{Research checkpoint, 23 November 2026}
\date{}
\begin{document}
\maketitle

This note gives a counterexample to one proposed transfer into attention
heads. It proves no lower bound for unrestricted head complexity. The
counterexample has only one input bit, so its target Boolean function itself
has a trivial one-head representation.

For $x\in[0,1]$, define the positive affine functions
\[
 a_1(x)=100-\frac{9999}{100}x,
 \qquad
 a_2(x)=a_3(x)=\frac19+\frac{80}{9}x.
\]
Let
\[
 P_t(x)=\prod_{h=1}^3\bigl(1+t(a_h(x)-1)\bigr),
 \qquad
 G_t(x)=\frac{\partial_t P_t(x)}{P_t(x)}
 =\sum_{h=1}^3\frac{a_h(x)-1}{1+t(a_h(x)-1)}.
\]
Every factor is positive for $0\leq t,x\leq1$, since it is a convex
combination of $1$ and a positive affine function. For $0<t\leq1$,
each of the three displayed denominators has strictly positive or strictly
negative slope in the original bit, and is positive throughout the interval.
Thus each summand is an admissible affine-ratio head under the repository's
positive- and negative-orientation realizations. At $t=0$ the fractions
reduce to affine functions; the argument below does not rely on treating
constant denominators as exact heads.

\begin{proposition}
The comparison $P_1(x)-P_0(x)$ is strictly positive at $x=0$ and
strictly negative at $x=1$. There is no common $t\in[0,1]$ for which
$G_t$ has these two strict signs.
\end{proposition}
\begin{proof}
At the two Boolean inputs,
\[
 (a_1(0),a_2(0),a_3(0))=(100,1/9,1/9),\qquad
 (a_1(1),a_2(1),a_3(1))=(1/100,9,9).
\]
Consequently
\[
 P_1(0)=100/81>1=P_0(0),\qquad
 P_1(1)=81/100<1=P_0(1).
\]
Subtracting fractions over positive denominators gives
\[
 G_t(0)=\frac{99}{1+99t}-\frac{16}{9-8t}
       =\frac{875-2376t}{(1+99t)(9-8t)},
\]
\[
 G_t(1)=-\frac{99}{100-99t}+\frac{16}{1+8t}
       =\frac{1501-2376t}{(100-99t)(1+8t)}.
\]
The denominators are positive on $[0,1]$. Hence $G_t(0)>0$ requires
$t<875/2376$, whereas $G_t(1)<0$ requires $t>1501/2376$.
These inequalities are incompatible.
\end{proof}

Pointwise, the fundamental theorem of calculus still gives
\[
 \log P_1(x)-\log P_0(x)=\int_0^1G_t(x)\,dt.
\]
The example shows that the pointwise mean-value parameter cannot be
chosen independently of the input. Therefore comparison of two products
of positive affine forms does not automatically compile to the logarithmic
derivative of their straight-line factor homotopy at one common parameter.
It does not exclude another construction, a quadrature with several
parameters, or a direct sign representation with different denominators.

All formulas above are a written proof. The two rational transition points
and their evaluations were additionally checked with Python's exact
\texttt{fractions.Fraction} arithmetic. No Lean theorem is asserted here.
\end{document}
